% =====================================================================
%  Formative Usability Evaluation of an Indonesian Home-Services
%  Marketplace Prototype  --  IEEE conference format (IEEEtran), 6 pages
%  Condensed build of paper/paper-1-usability-evaluation.md
%  Build: pdflatex paper-1-usability-evaluation-ieee.tex   (x2)
%
%  Reference metadata: every DOI below was resolved against the Crossref
%  REST API on 2026-10-10; authors, title, journal, volume, issue,
%  pagination and DOI are transcribed from the returned record.
% =====================================================================
\documentclass[conference]{IEEEtran}

\usepackage[T1]{fontenc}
\usepackage[utf8]{inputenc}
\usepackage{times}
\usepackage{amsmath}
\usepackage{amssymb}
\usepackage{booktabs}
\usepackage{array}
\usepackage{graphicx}
\usepackage[hidelinks]{hyperref}

\newcommand{\todo}{\mbox{\textit{[\ldots]}}}

% tighter float/caption spacing to reach the 6-page target
\setlength{\textfloatsep}{7pt plus 2pt minus 2pt}
\setlength{\floatsep}{5pt plus 2pt minus 2pt}
\setlength{\intextsep}{6pt plus 2pt minus 2pt}
\setlength{\abovecaptionskip}{4pt}
\setlength{\belowcaptionskip}{0pt}

\begin{document}

\title{Formative Usability Evaluation of an Indonesian Home-Services\\
Marketplace Prototype: Task Success, a Ceiling Effect, and a\\
Six-Theme Barrier Taxonomy}

\author{
\IEEEauthorblockN{Meutia Braniwati, S.Ds., M.Sn.}
\IEEEauthorblockA{BINUS University}
\IEEEauthorblockA{\todo{} verify affiliation and postal address before submission}
\IEEEauthorblockA{Corresponding author: \todo{} insert institutional email}
}

\maketitle

\begin{abstract}
We report a formative usability evaluation of a high-fidelity,
Indonesian-language web prototype for on-demand household maintenance
services. Ten participants completed eight standardised tasks across two
moderated rounds, yielding eighty task attempts, Single Ease Question
(SEQ) ratings on a 7-point scale, and observer commentary. Task success was
high at 93.75\% (75/80) and mean ease was $M=6.56$, $SD=1.05$. We argue
that this aggregate is close to uninformative. Computing dispersion per
task shows that four of eight tasks have saturated ($SD\leq0.42$) and can
no longer discriminate between a good and a very good design, while one
task---locating the FAQ and contact channel---spans the entire scale
($M=5.70$, $SD=2.11$, range 1--7) and is also the lowest-scoring. A matched
pair of discovery tasks isolates a taxonomy defect from a layout defect: a
well-labelled category was found without difficulty ($M=6.80$, $SD=0.42$)
whereas two near-synonymous labels cost 0.60 points and tripled the
variance ($M=6.20$, $SD=0.92$). We synthesise six barrier themes with
per-corpus denominators and recommend that formative evaluation report
per-task dispersion rather than aggregate success. We publish nine threats
to validity, four of them severe, and withdraw six unverifiable sources
from an earlier draft.
\end{abstract}

\begin{IEEEkeywords}
usability testing, Single Ease Question, ceiling effect, formative
evaluation, task success, MSME services, Indonesia, think-aloud
\end{IEEEkeywords}

% =====================================================================
\section{Introduction}
\label{sec:intro}

\subsection{Background and Motivation}

Household maintenance services in Indonesia---air-conditioner cleaning,
plumbing leak repair, pump servicing, septic tank suction---are dominated
by micro and small enterprises whose customer acquisition still runs
through \emph{spanduk} (street banners), printed leaflets, and vehicle
stickers~\cite{tianti}. That channel has structural properties invisible
to both firm and customer: a fixed physical area, capping how much
information can be presented; no record of prior work; and no price, which
is negotiated privately and is therefore unverifiable in advance. A web
interface removes the area constraint---but it does not follow that the
resulting interface is good, nor that it is measurable.

Most formative evaluations report a headline success rate, which is a
lossy summary of a multidimensional, per-task, per-participant result.
This paper shows, on real data, \emph{how much it hides.} Across ten
participants and eighty task attempts we find that 93.75\% success
coexists with (a)~four saturated tasks that can no longer generate
evidence, (b)~one task whose ratings span the entire scale, and (c)~a
service taxonomy that misleads participants who are otherwise fluent. Our
recommendation follows directly: \emph{formative evaluation should report
per-task dispersion and defect-probe results, not just totals.} A companion
paper~\cite{braniwati} treats the price-integrity finding separately,
raising a distinct methodological question that this paper sets aside.

\subsection{Research Questions}

\textbf{RQ1.} With what ease do young adult users complete a standardised
service-discovery and booking sequence? \textbf{RQ2.} Which tasks
discriminate between participants, and which are saturated? \textbf{RQ3.}
What are the characteristic barriers, and how frequently does each occur?
\textbf{RQ4.} Where do participant observation and independent code
inspection converge, and where does each uniquely contribute?

\subsection{Related Work}
\label{sec:related}

\textbf{Interface design for Indonesian MSME services.} Task-based design
and evaluation of e-commerce interfaces is well established in the
Indonesian literature~\cite{dinata,saputra,zunaidi,zakiyamani,nugrahaini};
Tianti \emph{et al.}~\cite{tianti} demonstrate a small-business website
raising promotional reach for a catering MSME, and Sunarso
\emph{et al.}~\cite{sunarso} show that visual content raises attraction and
conversion in MSME digital marketing. Recent national evidence frames the
opportunity: e-commerce platform adoption is associated with MSME income
growth~\cite{roihan}, while digital management-system
adoption~\cite{nurachmad} and digital literacy~\cite{muhtadin} predict MSME
performance. We adopt the underlying premise---that consistent visual
identity and legible service information signal operational
capability---while declining to import effect sizes measured on other
stimuli.

\textbf{Cognitive principles as rationale, not as measured effects.} Our
design drew on Norman's three levels of action design and the associated
grouping and isolation principles~\cite{norman}. These function as
\emph{rationale}: none was manipulated as an independent variable, and we
claim no causal effect.

\textbf{Formative evaluation instruments and their limits.} SEQ is a
widely adopted post-task instrument---one 7-point item administered
immediately after each attempt~\cite{sauro}---recently applied both alone
and against SUS~\cite{khodijah,azizah,purnamasari}. A systematic review of
applied usability metrics concluded that instrument choice governs which
usability attributes are measurable at all~\cite{wronikowska}. We adopt
SEQ for breadth---one comparable number per attempt---while noting a
trade-off our own data exposes sharply (\S\ref{sec:ease}): SEQ is a
\emph{mean} measure, and means conceal exactly what a formative study
needs.

\textbf{Ceiling effects.} Item ceiling effects are a documented hazard of
bounded self-report scales: they compress variance, bias estimates of
change, and can leave a scale unable to separate adjacent levels of
quality~\cite{edwards,sabella}. We treat this hazard as the object of
study rather than a caveat. Relatedly, recent work distinguishes formative
studies, which diagnose, from summative studies, which measure against
thresholds~\cite{schnitzbauer,salwei}; this paper is formative throughout,
and no variable measuring conversion, revenue, bounce rate, cognitive load,
or trust change was instrumented.

Heuristic inspection~\cite{nielsen} and its Indonesian
applications~\cite{perdanakusuma,febriyanthi,rabila,frobenius} are
complementary but operate on a different object: the code, not the
experience. \S\ref{sec:convergence} compares what each method caught.

% =====================================================================
\section{Methods}
\label{sec:methods}

\subsection{Prototype and Participants}

The prototype is a high-fidelity React application localised into
Indonesian, implementing service discovery with search and category
filtering, service detail pages, a four-stage checkout (summary
$\rightarrow$ schedule $\rightarrow$ address $\rightarrow$ payment), a
confirmation/invoice view, technician tracking, and an FAQ/contact area.
Two properties were left unchanged because they are the subject of the
research questions: prices render in \textbf{USD} while the interface is
Indonesian-language, and service descriptions include
\textbf{technician-internal procedure content} intended for engineers
rather than customers.

Ten participants were recruited and tested in two rounds of five,
identified only by anonymous codes (P01--P05, P06--P10). \textbf{No source
record contains age, gender, occupation, or any other demographic
attribute.} The protocol defined participants as Gen~Z users aged 18--26,
but that criterion \textbf{cannot be verified from the data collected},
and one session record reports readability difficulty ``for the
40-and-over age range,'' which is inconsistent with it. We therefore
describe the sample as \emph{young adult users recruited purposively} and
make \textbf{no representativeness claim} for Indonesian Gen~Z.

\subsection{Tasks, Measures, and Procedure}

Eight tasks span the booking flow (Fig.~\ref{fig:sequence}): T1 locate the
AC cleaning service; T2 locate the leak-detection service; T3 read the
detailed price breakdown; T4 select a booking date and slot; T5 complete
the address form; T6 trigger and identify a validation error; T7 navigate
back without losing entered data; T8 locate FAQ/contact support. T1 and T2
are \emph{matched} discovery tasks differing only in service category, so
their contrast isolates taxonomy from layout; T6 and T7 are \emph{defect
probes} with explicit failure criteria. Measures were \textbf{task
outcome} against a written criterion per task; \textbf{SEQ}, 1~=~\emph{very
difficult} $\ldots$ 7~=~\emph{very easy}, immediately after each attempt;
and \textbf{observer commentary}. Sessions were moderated with concurrent
think-aloud. Recording evidently occurred, but \textbf{no recording files
are present in the repository}; all qualitative material is therefore
observer paraphrase, not verbatim transcript.

\begin{figure*}[!t]
\centering
\includegraphics[width=\textwidth]{figures/p1-fig1-task-sequence.pdf}
\caption{The eight-task instrument mapped onto the prototype's booking
sequence. Six sequential screens carry T1--T6; T7 and T8 probe behaviour
available from any point in the flow.}
\label{fig:sequence}
\end{figure*}

\subsection{Resolving a Completion-Rate Dispute}

Completion was coded during sessions and recorded as 39/40 (round~1) and
40/40 (round~2). Both source recap sheets then appended an auditor note
that specific attempts should have been coded failures: in round~1, two T6
attempts recorded as successes although the criterion was \emph{find the
validation message} and the notes state the participants did not find it;
in round~2, one T6 attempt recorded as success although the note reads
``user did not find any error,'' and another reached confirmation with a
street address consisting only of \texttt{(} and \texttt{)}. This is an
interpretive disagreement about whether failure on a \emph{sub-criterion}
constitutes task failure, not an arithmetic error. Rather than adopt
either position silently, \textbf{we report the auditor-corrected figures
as the headline and the as-recorded figures as a sensitivity analysis}
(\S\ref{sec:success}); the two differ by four attempts---larger than any
other coding uncertainty in this study.

% =====================================================================
\section{Results}
\label{sec:results}

\subsection{Task Success}
\label{sec:success}

\begin{table}[!t]
\caption{Task completion as recorded and as auditor-corrected.}
\label{tab:success}
\centering
\footnotesize
\setlength{\tabcolsep}{4pt}
\begin{tabular}{@{}lccc@{}}
\toprule
 & Round 1 ($n{=}5$) & Round 2 ($n{=}5$) & Combined ($n{=}10$) \\
\midrule
As recorded       & 39/40 (97.5\%) & 40/40 (100\%) & 79/80 (\textbf{98.75 \%}) \\
Auditor-corrected & 37/40 (92.5\%) & 38/40 (95.0\%) & 75/80 (\textbf{93.75 \%}) \\
 & (headline) & (headline) & (headline) \\
\bottomrule
\end{tabular}
\end{table}

All five failures occurred at T6 or T8. A rate above 90\% would ordinarily
indicate a largely successful design. \textbf{We do not read it that
way}: the T6 criterion was permissive; success is not integrity, since
participants reached confirmation while the interface's own figures were
internally inconsistent~\cite{braniwati}; and, decisively, success is
\emph{ceiling-bound}---four tasks were near-universally completed, so the
metric cannot discriminate (\S\ref{sec:ease}).

\subsection{Ease, and Why the Mean Is Not the Finding}
\label{sec:ease}

Overall SEQ: \textbf{$M=6.56$, $SD=1.05$} ($n=80$), median 7; distribution
62 ratings of 7, 10 of 6, 4 of 5, 1 of 4, 2 of 3, 1 of 1. Two observations
follow, and \textbf{neither is visible in the aggregate mean.}

\textbf{(a) Four of eight tasks are saturated.} T1, T3, T4, and T5 have
$SD\leq0.42$; for T4 and T5 no participant rated below 6. These tasks
cannot separate a good design from a very good one, and further iterations
must not rely on them for evidence. Reporting only the mean would have
hidden this entirely---and would have invited us to keep iterating on tasks
that had already stopped producing signal.

\textbf{(b) T8 is the study's real signal.} Support-channel
discoverability has the lowest mean (\emph{M}~=~5.70) \emph{and} by far
the widest dispersion (SD~=~2.11), spanning the full usable range. The
comments locate the cause precisely: participants found the WhatsApp
control and read it as a \emph{booking} action rather than a
\emph{support} channel; some located help only in the footer; one found no
direct contact number at all. Two rated T8 at 7; one rated it 1.
\textbf{A task whose ratings span the entire scale is under-served by the
design, not merely unfamiliar to the participant.} A wide spread invites
the reading ``some users need training''; the comments refute it. The
mechanism is not comprehension but \emph{labelling and placement}---a
control that is present, visible, and means something other than it
appears to.

\begin{figure*}[!t]
\centering
\includegraphics[width=\textwidth]{figures/p1-fig2-seq-ceiling.pdf}
\caption{Mean SEQ per task with $\pm 1$ SD whiskers. Four tasks sit against
the scale ceiling with $SD\leq0.42$. T8 has both the lowest mean and the
widest dispersion, spanning the full 1--7 range.}
\label{fig:ceiling}
\end{figure*}

\subsection{The Matched Pair Isolates a Taxonomy Defect}
\label{sec:matched}

T1 and T2 are the same task with a different category. T1 (find AC
cleaning) reached \textbf{$M=6.80$, $SD=0.42$}; T2 (find leak detection)
reached \textbf{$M=6.20$, $SD=0.92$}---a 0.60-point drop with triple the
variance, on an otherwise identical task. Participants attributed the
difficulty to near-synonymous service labels (``smart leak detection''
versus pipe inspection) that the interface did not differentiate, and in
two cases to the booking control sitting below the fold. This is a
\textbf{labelling} problem, not a \textbf{layout} problem: the category was
reachable, but participants could not determine \emph{which} entry
addressed their problem. The practical consequence is disproportionate,
since merging two catalogue entries is a data change while redesigning the
discovery screen is a redesign.

\subsection{Barrier Taxonomy}
\label{sec:themes}

Six themes emerged. Because the narrative corpus and the SEQ corpus are
not confirmed to be the same people (\S\ref{sec:threats}, T-1), counts are
reported with explicit denominators rather than pooled.

\begin{table*}[!t]
\caption{Six-theme barrier taxonomy with per-corpus denominators. Counts
are not pooled across corpora because the two rosters are not confirmed to
be the same participants (\S\ref{sec:threats}, T-1).}
\label{tab:themes}
\centering
\scriptsize
\begin{tabular}{@{}>{\raggedright\arraybackslash}p{0.9cm}>{\raggedright\arraybackslash}p{2.6cm}>{\raggedright\arraybackslash}p{2.2cm}>{\raggedright\arraybackslash}p{1.8cm}>{\raggedright\arraybackslash}p{7.1cm}@{}}
\toprule
\textbf{\#} & \textbf{Theme} & \textbf{Narrative} & \textbf{SEQ} &
\textbf{Observed barrier} \\
 & & \textbf{(of 10)} & \textbf{(of 10)} & \\
\midrule
\textbf{T1} & Currency \& technical register & 3 named (synthesis asserts
  unanimity) & 2---both round 2, \textbf{none in round 1} & USD on an
  Indonesian site; jargon (``precision diagnostic'', ``internal
  evaporator'', ``indoor--outdoor unit relation'') over-excludes lay
  users; technician-internal documentation shown to customers \\
\textbf{T2} & Price transparency & 3 & 1 (+1 currency remark) & Price
  breakdown not collapsible; unit scope (which AC, indoor or outdoor)
  unstated \\
\textbf{T3} & Scheduling rigidity & 6 & 5 & Date selected twice with
  \textbf{disagreeing} options; 3--5 day window read as weekday-only; no
  calendar; no arrival range \\
\textbf{T4} & Address capture \& validation & 7 & 5 & No map, geocoding, or
  autocomplete; no saved address; postal code validated but street address
  not; mandatory fields unmarked; ``service address'' misread as the
  provider's office \\
\textbf{T5} & Navigation \& state persistence & 6 & 8 & Form data lost on
  Back/header navigation and on reload; footer links return to the home
  page; search query persists after clearing; confirmation button stuck in
  ``processing'' \\
\textbf{T6a/b/c} & Support, confirmation clutter, accessibility & 4 / 1 / 2
  & 4 / 1 / 1 & Help buried in footer; three-item FAQ; duplicate WhatsApp
  and print buttons, two tracking bars, unexplained QR code; type size and
  grey-on-white contrast below threshold \\
\bottomrule
\end{tabular}
\end{table*}

\textbf{One finding narrows rather than repeats an earlier complaint.} On
scheduling (T3), a respondent stated that seeing the date \emph{twice} is
acceptable if it functions as a cross-check---\textbf{provided the two
instances agree}. The defect is therefore \textbf{disagreement between two
date pickers}, not duplication. A single source of truth is consequently a
smaller change than removing one picker, and a reviewer who treats
duplication as the problem would prescribe the wrong fix.

\subsection{Convergence with Code Inspection}
\label{sec:convergence}

An independent code-based heuristic audit produced 0 Critical, 2 Major, and
4 Minor findings. \textbf{Four of six audited defects were independently
rediscovered by participants without prompting}, which supports heuristic
inspection as a cheap pre-test filter (Table~\ref{tab:convergence}). The
two misses---USD on a localised site, and a total inconsistent with its
own line items---motivate the companion paper~\cite{braniwati}. Because
the audit records these defects as \emph{repaired} while the tested
sessions reported all of them, the tested build is probably
pre-remediation (\S\ref{sec:threats}, T-2).

\begin{table}[!t]
\caption{Convergence between code-audit findings and participant-observed
barriers.}
\label{tab:convergence}
\centering
\scriptsize
\begin{tabular}{@{}>{\raggedright\arraybackslash}p{1.5cm}>{\raggedright\arraybackslash}p{6.6cm}@{}}
\toprule
\textbf{Verdict} & \textbf{Code-audit finding $\rightarrow$ observed barrier} \\
\midrule
\textbf{Converges} & Major: home-page price differs from catalogue/detail
  $\rightarrow$ T2 unexplained total jumps \\
\textbf{Not testable} & Major: no visible keyboard focus indicator
  $\rightarrow$ --- \\
\textbf{Converges} & Minor: header section links fragile $\rightarrow$ T5
  footer links returned users to home \\
\textbf{Converges} & Minor: stepper icons contradict text $\rightarrow$ T6b
  duplicate tracking status bars \\
\textbf{Converges} & Minor: no re-order path $\rightarrow$ T5 participants
  requested order history \\
\textbf{Partly} & Minor: postal code silently truncated $\rightarrow$ T4
  postal code validated, street address not \\
\textbf{Missed} & --- $\rightarrow$ T1 USD currency on a localised site \\
\textbf{Missed} & --- $\rightarrow$ T2 total $\neq$ sum of its own line
  items \\
\bottomrule
\end{tabular}
\end{table}

\subsection{A Contradiction Between Our Own Sources}

The narrative report's recommendation matrix rates support/contact
remediation \textbf{Low} priority; both recap sheets rate it \textbf{High},
and four of the narrative's own sessions record difficulty locating
support. We adopt the recap sheets' rating, on the grounds that it is
supported by the largest weight of recorded evidence and by the dispersion
analysis in \S\ref{sec:ease}(b). We report the disagreement rather than
suppress it, so that a reader checking our recommendation against the
primary documents can see that we noticed.

% =====================================================================
\section{Discussion}
\label{sec:discussion}

\subsection{What Worked}

The prototype's \emph{structure} is sound. Four tasks sit at or near
ceiling with almost no dispersion, the booking flow follows an expected
mental model---browse, detail, schedule, address, pay, confirm---and
participants navigated it without instruction. Note the asymmetry in T5:
the address form reached ceiling (\emph{M}~=~6.90, SD~=~0.32)
\textbf{while generating the most frequently requested features in the
entire study}---geocoding, autocomplete, saved addresses. The form is easy
to fill and wrong in design. \textbf{Ease of completion and fitness for
purpose are independent properties, and SEQ measures only the first.}

Three patterns organise the failures. \textbf{Pattern 1---presentation
locale not matched to content locale:} an Indonesian-language interface
renders prices in dollars and service descriptions in technician-level
English. No individual string is wrong; the locale as a whole is
incoherent, and participants read it as evidence the product was not built
for them. \textbf{Pattern 2---missing infrastructure treated as an
acceptable default:} no map, no autocomplete, no saved address, no order
history, no calendar---collectively, a checkout that asks the user to do
work the platform could do, and to do it again next booking.
\textbf{Pattern 3---redundant, conflicting surfaces:} two date pickers
that disagree, a home-page price differing from the catalogue, two
tracking bars and an unexplained QR code. Redundancy is normally benign;
here it is actively harmful, because when two copies of a fact disagree
the user cannot tell which to trust---and in a pricing context, they will
trust neither.

\begin{table}[!t]
\caption{Prioritised remediation backlog. Each row gives the finding and
its recommended action.}
\label{tab:backlog}
\centering
\scriptsize
\setlength{\tabcolsep}{2.4pt}
\begin{tabular}{@{}c >{\raggedright\arraybackslash}p{5.9cm} >{\raggedright\arraybackslash}p{1.25cm}@{}}
\toprule
\textbf{P} & \textbf{Finding $\rightarrow$ recommended action} & \textbf{Sev} \\
\midrule
1 & Total inconsistent with its own line items $\rightarrow$ add a
  build-time invariant asserting \texttt{total == service + fee + tax}
  & \textbf{Critical} \\
2 & USD on an Indonesian site; technician jargon $\rightarrow$ localise to
  IDR; rewrite the taxonomy in lay register; merge near-synonymous
  entries & \textbf{High} \\
3 & Support channel undiscoverable $\rightarrow$ persistent, explicitly
  labelled support action in the header; direct contact; FAQ search
  & \textbf{High} \\
4 & Form state lost on Back/reload $\rightarrow$ persist checkout state;
  warn before destructive navigation; offer resume & \textbf{High} \\
5 & Street address unvalidated; required fields unmarked $\rightarrow$
  validate street address as postal code is; mark required fields; relabel
  ``service address'' & \textbf{High} \\
6 & Two date pickers that disagree $\rightarrow$ one source of truth for
  available dates, with an identical read-only cross-check
  & \textbf{High} \\
7 & Footer links return to home; query persists $\rightarrow$ fix footer
  routing; reset the query on clear & \textbf{Medium} \\
8 & Confirmation-page clutter $\rightarrow$ collapse to a static summary;
  move tracking to order history; label or remove the QR code
  & \textbf{Medium} \\
9 & Type size and contrast below threshold $\rightarrow$ minimum 12\,pt;
  secondary text to WCAG AA; verify rendered contrast & \textbf{Medium} \\
10 & No geocoding, saved address, order history $\rightarrow$ map pin and
  autocomplete; addresses stored by label; order history with re-order
  & \textbf{Medium} \\
11 & No path for a missing service; rigid scheduling $\rightarrow$ ``other
  service'' request; full calendar; arrival-time range & \textbf{Low} \\
\bottomrule
\end{tabular}
\end{table}

\subsection{Methodological Recommendations}

\begin{itemize}
\item \textbf{Report per-task dispersion, not just aggregate success.} A
93.75\% success rate coexisted here with a Critical pricing defect and a
task rated as low as 1/7.
\item \textbf{Separate discriminating tasks from saturated ones.} Four of
our eight tasks could no longer produce evidence; iterating against them
consumes budget for no return.
\item \textbf{Prefer matched task pairs to isolate causes.} The T1/T2
contrast converted a vague complaint (``search is hard'') into a specific,
cheap-to-fix diagnosis.
\item \textbf{Treat effortful completion as a smell.} The address form was
easy to fill and wrong. Comfort is not fitness.
\end{itemize}

\subsection{Threats to Validity}
\label{sec:threats}

\noindent\textbf{T-1---Participant rosters disagree across source records.
[Severe]} The narrative corpus names ten participants; the two recap
sheets name a partly different set---four overlap, three appear only in the
recap sheets, six narrative participants have no SEQ record anywhere.
\textbf{Consequence:} the corpora cannot be confirmed to be the same ten
people. We never aggregate across them as though they were, and every
count in \S\ref{sec:themes} carries an explicit denominator.

\medskip\noindent\textbf{T-2---No build identifier, date, or moderator record.
[Severe]} The protocol required recording the canonical prototype and its
build hash before sessions began; no session record contains one, and the
tested build probably predates the recorded audit remediation.
\textbf{Consequence:} no result may be attributed to the current
prototype.

\medskip\noindent\textbf{T-3---Demographics unrecorded; Gen~Z eligibility
unverifiable. [Severe]} \textbf{Consequence:} for a paper framed around
young users, the sample's composition is unestablished.

\medskip\noindent\textbf{T-4---Ethics and consent unevidenced. [Severe]}
No consent record, approval reference, or clearance number appears in any
source report. \textbf{Consequence:} an unresolved compliance gap that
must be closed before publication of this data.

\medskip\noindent\textbf{T-5---Ceiling effect limits four tasks.
[Moderate]} T1, T3, T4, T5 ($SD\leq0.42$) cannot discriminate between
design states; later iterations will need harder tasks or a more complex
stimulus.

\medskip\noindent\textbf{T-6---Variance collapse between rounds.
[Moderate]} Round-2 per-task SD (0.49) is $2.8\times$ smaller than
round-1 (1.39), with three participants returning near-uniform maximum
scores. Differing scoring discipline, moderator effects, demand
characteristics, and transcription convention cannot be distinguished from
the available records, so we draw \textbf{no} design-effect conclusion.

\medskip\noindent\textbf{T-7---Qualitative corpus is observer paraphrase,
not transcript. [Moderate]} \textbf{Consequence:} no direct quotation is
used anywhere in this paper.

\medskip\noindent\textbf{T-8---Transcription gaps. [Low]} One recap sheet
notes that a participant's comments were truncated in a screenshot; the
check against the original sheet was not performed.

\medskip\noindent\textbf{T-9---References partly unverified. [Low]} Two
citations remain only partially verified (see the note after the reference
list); six sources from an earlier draft were withdrawn.

% =====================================================================
\section{Conclusion}

We evaluated an Indonesian-language home-services marketplace prototype
with ten participants across eight tasks and two rounds. Task success was
93.75\% and mean ease 6.56/7---and the most useful result is that
\textbf{neither number characterises the design.} Beneath the aggregate sit
four saturated tasks that can no longer generate evidence, one task
spanning the entire scale because a help channel is present but
mislabelled, and a matched pair isolating a taxonomy defect from a layout
defect. The failures are concentrated and diagnosable: a locale that is
Indonesian in language and American in currency, a checkout that discards
state on navigation, and two date pickers that disagree with each other.

The lesson generalises beyond this prototype: a formative study reporting
only aggregate success rates will systematically under-report exactly the
problems that matter most, because discriminating tasks are outnumbered by
saturated ones. We recommend publishing per-task dispersion as a matter of
course, and we have published ours---with the nine threats to validity that
constrain them.

% =====================================================================
\begin{thebibliography}{}
\scriptsize

\bibitem{roihan}
R.~I.~Roihan, B.~Yateno, and J.~Imantoro, ``E-commerce platform adoption
and MSME income growth: Evidence from Indonesia,'' \emph{Krakatoa Mgmt.
Res. J.}, vol.~1, no.~1, pp.~1--13, 2026. doi:
\href{https://doi.org/10.61401/kmrj.v1i1.544}{10.61401/kmrj.v1i1.544}

\bibitem{nurachmad}
E.~Nurachmad, M.~Mulyana, A.~Sukamto, and I.~Mulia, ``Adoption of
digital-based management information systems and its impact on MSME
performance in Indonesia,'' \emph{J. Ilm. Manaj. Kesatuan}, vol.~13, no.~1,
pp.~629--634, 2025. doi:
\href{https://doi.org/10.37641/jimkes.v13i1.4004}{10.37641/jimkes.v13i1.4004}

\bibitem{muhtadin}
I.~Muhtadin \emph{et al.}, ``The influence of digital entrepreneurship
and digital literacy on MSME performance through digital technology
adoption in Indonesia,'' \emph{The Es Economics and Entrepreneurship},
vol.~4, no.~3, 2026. doi:
\href{https://doi.org/10.58812/esee.v4i03.978}{10.58812/esee.v4i03.978}

\bibitem{dinata}
P.~Z.~Dinata, M.~A.~Urwah, M.~R.~Rahmawan, and E.~Junaeti,
``Perancangan UI/UX web e-commerce \emph{Hallo Coffee} menggunakan metode
\emph{user centered design},'' \emph{Jambura J. Inform.}, vol.~5, no.~1,
pp.~45--58, 2023. doi: \href{https://doi.org/10.37905/jji.v5i1.17511}{10.37905/jji.v5i1.17511}

\bibitem{sunarso}
B.~Sunarso, Tusriyanto, and F.~Mustafa, ``Analysing the role of visual
content in increasing attraction and conversion in MSME digital
marketing,'' \emph{J. Contemp. Admin. Manag. ADMAN}, vol.~1, no.~3,
pp.~193--200, 2023. doi: \href{https://doi.org/10.61100/adman.v1i3.79}{10.61100/adman.v1i3.79}

\bibitem{tianti}
A.~M.~Tianti \emph{et al.}, ``Peningkatan promosi usaha catering UMKM melalui
pembuatan website berbasis WordPress,'' \emph{PENITI BANGSA}, vol.~2,
no.~2, 2024. doi: \href{https://doi.org/10.30742/peniti-bangsa.v2i22024.177}{10.30742/peniti-bangsa.v2i22024.177}

\bibitem{saputra}
A.~D.~Saputra, ``Usability evaluation of the Alodokter application using
the task-based usability testing method,'' \emph{J. Tek. Ind.}, vol.~11,
no.~2, pp.~294--302, 2025. doi:
\href{https://doi.org/10.24014/jti.v11i2.38790}{10.24014/jti.v11i2.38790}

\bibitem{zunaidi}
R.~A.~Zunaidi, A.~D.~R.~Purbantari, and H.~Mardhiana, ``Usability testing
of the online platform for selling foodstuffs in Indonesia,''
\emph{J. Ind. Eng. Manag.}, vol.~6, no.~3, pp.~18--28, 2021. doi:
\href{https://doi.org/10.33536/jiem.v6i3.948}{10.33536/jiem.v6i3.948}

\bibitem{zakiyamani}
M.~Zakiyamani and L.~P.~Manik, ``Usability aplikasi kebencanaan di
Indonesia dengan usability testing dan sistem usability scale,''
\emph{INTECOMS}, vol.~5, no.~2, pp.~272--281, 2023. doi:
\href{https://doi.org/10.31539/intecoms.v5i2.5335}{10.31539/intecoms.v5i2.5335}

\bibitem{nugrahaini}
K.~N.~Nugrahaini and E.~Nurmiati, ``Usability evaluation of social
commerce integration on TikTok Shop among generation Z using system
usability scale,'' \emph{J. Inform. Rekayasa Perangkat Lunak}, vol.~7,
no.~3, pp.~168--176, 2026. doi:
\href{https://doi.org/10.33365/jatika.v7i3.1850}{10.33365/jatika.v7i3.1850}

\bibitem{khodijah}
K.~Khodijah, R.~M.~Febriyanti, and A.~Y.~Rindarwati, ``Evaluasi penggunaan
aplikasi diabetes M-Blood Sugar Diary menggunakan metode Single Ease
Question (SEQ) dan System Usability Scale (SUS),'' \emph{J. Penelit. Sains},
vol.~26, no.~2, art.~no.~979, 2024. doi:
\href{https://doi.org/10.56064/jps.v26i2.979}{10.56064/jps.v26i2.979}

\bibitem{azizah}
N.~A.~Azizah, I.~N.~T.~A.~Putra, A.~Bendesa, S.~Abdullah, and
G.~Y.~Prakoso, ``Analisis kebutuhan UI/UX aplikasi mobile travel planner
Wonderworld dengan design thinking dan Single Ease Question (SEQ),''
\emph{J. Inform. Tek. Elektro Terapan}, vol.~13, no.~2, 2025. doi:
\href{https://doi.org/10.23960/jitet.v13i2.6460}{10.23960/jitet.v13i2.6460}

\bibitem{purnamasari}
S.~A.~Purnamasari, N.~Heryana, and K.~Prihandani, ``Perbandingan
penggunaan System Usability Scale dan Usefull, Satisfaction and Ease of
Use Questionnaire pada usability testing,'' \emph{J. Ilm. Inform.}, vol.~6,
no.~1, pp.~59--69, 2021. doi:
\href{https://doi.org/10.35316/jimi.v6i1.1236}{10.35316/jimi.v6i1.1236}

\bibitem{wronikowska}
M.~W.~Wronikowska \emph{et al.}, ``Systematic review of applied usability
metrics within usability evaluation methods for hospital electronic
healthcare record systems,'' \emph{J. Eval. Clin. Pract.}, vol.~27, no.~6,
pp.~1403--1416, 2021. doi:
\href{https://doi.org/10.1111/jep.13582}{10.1111/jep.13582}

\bibitem{edwards}
K.~D.~Edwards and J.~Soland, ``How scoring approaches impact estimates of
growth in the presence of survey item ceiling effects,'' \emph{Appl.
Psychol. Meas.}, vol.~48, no.~3, pp.~147--164, 2024. doi:
\href{https://doi.org/10.1177/01466216241238749}{10.1177/01466216241238749}

\bibitem{sabella}
S.~A.~Sabella and S.~T.~Meier, ``Pervasive ceiling effect in supervisory
working alliance measurement: An exploratory meta-analysis,''
\emph{Clin. Supervisor}, vol.~43, no.~2, pp.~248--268, 2024. doi:
\href{https://doi.org/10.1080/07325223.2024.2393299}{10.1080/07325223.2024.2393299}

\bibitem{schnitzbauer}
A.~A.~Schnitzbauer \emph{et al.}, ``Usability testing of a digitized
interventional prehabilitation tool for health care professionals and
patients before major surgeries: Formative and summative evaluation,''
\emph{JMIR Formative Res.}, vol.~8, art.~no.~e59513, 2024. doi:
\href{https://doi.org/10.2196/59513}{10.2196/59513}

\bibitem{salwei}
M.~E.~Salwei \emph{et al.}, ``Evaluating the safety and usability of an
over-the-counter medical device for adults with mild to moderate hearing
loss: Formative and summative usability testing,'' \emph{JMIR Human
Factors}, vol.~12, art.~no.~e65142, 2025. doi:
\href{https://doi.org/10.2196/65142}{10.2196/65142}

\bibitem{nielsen}
J.~Nielsen, ``Enhancing the explanatory power of usability heuristics,'' in
\emph{Proc. ACM CHI'94}, 1994, pp.~152--158. doi:
\href{https://doi.org/10.1145/191666.191729}{10.1145/191666.191729}

\bibitem{perdanakusuma}
A.~R.~Perdanakusuma, B.~T.~Hanggara, and A.~R.~Hasnanursanti, ``Analisis
usability website resmi pemerintah kota \emph{Surakarta} menggunakan
metode heuristic evaluation,'' \emph{J. TECNOSCIENZA}, vol.~6, no.~2,
pp.~429--443, 2022. doi:
\href{https://doi.org/10.51158/tecnoscienza.v6i2.736}{10.51158/tecnoscienza.v6i2.736}

\bibitem{febriyanthi}
A.~Febriyanthi, ``Interface design heuristic evaluation website yayasan
kesehatan Telkom,'' \emph{Turk. J. Comput. Math. Educ. TURCOMAT}, vol.~12,
no.~4, pp.~852--859, 2021. doi:
\href{https://doi.org/10.17762/turcomat.v12i4.573}{10.17762/turcomat.v12i4.573}

\bibitem{rabila}
S.~Rabila, S.~Munir, T.~Nabarian, and D.~Wismanindra, ``Analisis dan
evaluasi user interface design untuk usability menggunakan metode
heuristic evaluation pada web perusahaan bioteknologi,'' \emph{J.
Informatika Terpadu}, vol.~8, no.~2, pp.~68--77, 2022. doi:
\href{https://doi.org/10.54914/jit.v8i2.420}{10.54914/jit.v8i2.420}

\bibitem{frobenius}
A.~C.~Frobenius and R.~C.~Kurniawan, ``User interface evaluation of the
Sumber Alam Ekspres application using the heuristic evaluation method,''
\emph{J. Appl. Inform. Comput.}, vol.~9, no.~3, pp.~651--657, 2025. doi:
\href{https://doi.org/10.30871/jaic.v9i3.9285}{10.30871/jaic.v9i3.9285}

\bibitem{norman}
D.~A.~Norman, \emph{The Design of Everyday Things: Revised and Expanded
Edition}. New York, NY, USA: Basic Books, 2013.

\bibitem{sauro}
J.~Sauro and J.~R.~Lewis, \emph{Quantifying the User Experience: Practical
Statistics for User Research}, 2nd~ed. Sebastopol, CA, USA: Morgan
Kaufmann, 2016.

\bibitem{braniwati}
M.~Braniwati, ``Formative usability evaluation of an Indonesian
home-services marketplace prototype,'' companion paper, same study, 2026.
\end{thebibliography}

\noindent\scriptsize\textit{Reference verification.} Every reference
above was resolved against the Crossref REST API on 2026-10-10; authors,
title, journal, volume, issue, pagination and DOI are transcribed from the
returned record. Two remain only partially verified and must be checked by
hand: \cite{perdanakusuma} (irregular deposited author forms; publisher
page range disagrees with earlier drafts) and \cite{febriyanthi} (author
field deposited as malformed ``Et. al.'').

% =====================================================================
\appendices

\section{Instrument, SEQ Matrix, and Provenance}
\label{app:seq}

\textbf{Instrument.} A single item, \emph{``After completing the task, how
easy or difficult was that task for you?''}, anchored 1~=~very difficult
and 7~=~very easy, administered immediately after each of T1--T8. Ordinal
scale; mean and SD are reported for comparability, with the caveat in
\S\ref{sec:threats}(T-5). Table~\ref{tab:seq} reproduces the matrix exactly
as recorded in the source recap sheets; recomputation confirmed every
reported mean.

\begin{table}[!t]
\centering
\caption{Complete Single Ease Question matrix.}
\label{tab:seq}
\scriptsize
\setlength{\tabcolsep}{2.6pt}
\begin{tabular}{@{}c|rrrrrrrr|r@{}}
\toprule
\textbf{P} & \textbf{T1} & \textbf{T2} & \textbf{T3} & \textbf{T4} &
\textbf{T5} & \textbf{T6} & \textbf{T7} & \textbf{T8} & \textbf{Mean} \\
\midrule
P01 & 6 & 5 & 7 & 7 & 7 & 7 & 4 & 6 & 6.13 \\
P02 & 7 & 7 & 7 & 7 & 7 & 7 & 7 & 7 & 7.00 \\
P03 & 7 & 5 & 7 & 7 & 7 & 7 & 7 & 3 & 6.25 \\
P04 & 7 & 7 & 7 & 7 & 7 & 7 & 7 & 7 & 7.00 \\
P05 & 7 & 5 & 7 & 7 & 7 & 3 & 7 & 1 & 5.50 \\
P06 & 6 & 6 & 6 & 6 & 6 & 7 & 7 & 5 & 6.13 \\
P07 & 7 & 6 & 7 & 7 & 7 & 7 & 7 & 7 & 6.88 \\
P08 & 7 & 7 & 7 & 7 & 7 & 7 & 7 & 7 & 7.00 \\
P09 & 7 & 7 & 7 & 7 & 7 & 7 & 7 & 7 & 7.00 \\
P10 & 7 & 7 & 6 & 7 & 7 & 6 & 7 & 7 & 6.75 \\
\midrule
\textbf{M} & \textbf{6.80} & \textbf{6.20} & \textbf{6.80} &
\textbf{6.90} & \textbf{6.90} & \textbf{6.50} & \textbf{6.70} &
\textbf{5.70} & \textbf{6.56} \\
\textbf{SD} & 0.42 & 0.92 & 0.42 & 0.32 & 0.32 & 1.27 & 0.95 & 2.11 &
1.05 \\
\bottomrule
\end{tabular}
\end{table}

\textbf{Provenance.} The full traceability record is maintained in the
study's integration record. Six sources cited in the earlier proposal could
not be located in Crossref, DataCite, the DOI Handle system, Google
Scholar, GARUDA, or SINTA, though several used plausible real author names
and real journal titles; all have been removed from the numbered list and
the body text, with dependent claims withdrawn, and the one apparently
fabricated entry (Hart \& Kapp, \emph{J. Des. Res. Pract.}, 2021) has no
trace of the article, the authors, or the journal anywhere. \textbf{Do not
re-add them without retrieving the actual record.} \textbf{Barred figures:}
the superseded draft's 87.5\%, 4.20/5, 84.6\%, and 4.23/5 were never
derivable from any source; 84.6\% is not the corrected completion rate, and
87.5\% appears nowhere in \S\ref{sec:success}.

\end{document}